\section{Related Work}

\subsection{The Building Energy Performance Gap}
The building energy performance gap refers to the systemic difference between the energy consumption predicted during the design phase of a building and the actual energy used once the building is operational \cite{performance_gap_1, performance_gap_2}. This gap can be substantial, with actual consumption frequently exceeding design predictions by a factor of two to three \cite{performance_gap_3}. The root causes are multifaceted, encompassing discrepancies in building physics modeling, construction quality, and, critically, occupant behavior and operational control strategies \cite{performance_gap_1}.

Addressing this gap requires moving beyond static, design-phase simulations toward continuous, data-driven operational auditing. However, traditional energy audits are labor-intensive, relying on manual inspections and the temporary installation of sub-metering equipment. While whole-building smart meters have become ubiquitous, providing high-resolution aggregate consumption data, they lack the granularity required to pinpoint inefficiencies in specific subsystems, such as lighting or HVAC. This limitation has spurred significant interest in computational methods capable of extracting subsystem-level insights from aggregate meter data.

\subsection{Non-Intrusive Load Monitoring (NILM)}
Non-Intrusive Load Monitoring (NILM), first conceptualized by Hart in the early 1990s \cite{hart_1992}, aims to disaggregate a total electrical load into its constituent appliance-level components using only a single measurement point. NILM is fundamentally a single-channel blind source separation problem.

Historically, NILM research has focused heavily on the residential sector, driven by the availability of high-frequency datasets (e.g., REDD \cite{redd_2011}, UK-DALE \cite{ukdale_2015}) that capture voltage and current waveforms at kHz or MHz sampling rates. At these frequencies, individual appliances exhibit unique, transient electrical signatures (e.g., the specific current draw profile of a refrigerator compressor starting) \cite{nilm_survey_1, nilm_survey_2}.

In contrast, commercial buildings present a significantly more challenging environment for NILM. Commercial smart meters typically record data at much lower frequencies (e.g., 15-minute or hourly intervals), completely obliterating transient signatures. Furthermore, commercial loads are characterized by high concurrency, where numerous devices operate simultaneously, creating a complex, overlapping aggregate signal \cite{nilm_survey_1}. Consequently, low-frequency commercial NILM must rely on macroscopic temporal patterns, such as diurnal cycles and operational schedules, rather than electrical transients.

\subsection{Deep Learning for Energy Disaggregation}
Early approaches to low-frequency disaggregation relied on classical time series decomposition techniques, such as Seasonal-Trend Decomposition using LOESS (STL) \cite{stl_1990}, or probabilistic models like Hidden Markov Models (HMMs) and their variants (e.g., Factorial HMMs) \cite{fhmm_2012}. While STL is effective for simple, periodic signals, it struggles to separate the complex, non-stationary dynamics inherent in commercial building loads. Similarly, FHMMs suffer from severe scalability issues due to the exponential growth of the state space as the number of appliances increases.

The advent of deep learning has revolutionized NILM. Kelly and Knottenbelt \cite{kelly_2015} were among the first to apply deep neural networks to energy disaggregation, demonstrating significant improvements over HMM-based methods. Subsequent research has explored a wide array of architectures. Zhang et al. \cite{zhang_2018} introduced a highly influential sequence-to-point (seq2point) learning framework using Convolutional Neural Networks (CNNs), which maps a window of aggregate power readings to the midpoint consumption of a target appliance.

More recently, attention mechanisms and Transformer architectures have been adapted for NILM. BERT4NILM \cite{bert4nilm} leverages a bidirectional Transformer encoder to capture long-range temporal dependencies across the aggregate signal, achieving state-of-the-art performance on several residential benchmarks. However, a critical limitation of these advanced deep learning models is their reliance on supervised learning. They require extensive datasets containing both the aggregate signal and the corresponding ground-truth sub-metered data for each target appliance.

\subsection{Unsupervised Time Series Analysis}
In the commercial sector, ground-truth sub-metered data is rarely available, necessitating unsupervised approaches. Unsupervised deep learning for time series often involves representation learning followed by clustering or anomaly detection. Convolutional Autoencoders (CAEs) have proven highly effective at extracting robust, shift-invariant features from time series data \cite{cae_2021}. By compressing the input sequence into a lower-dimensional latent space, CAEs force the network to learn the most salient structural patterns, filtering out high-frequency noise.

For the disaggregation task itself, our work draws inspiration from the N-BEATS (Neural Basis Expansion Analysis for Interpretable Time Series) architecture \cite{nbeats_2020}. Originally proposed by Oreshkin et al. for univariate time series forecasting, N-BEATS employs a deep stack of fully connected layers with forward and backward residual links. A key innovation of N-BEATS is its interpretability: the network can be constrained to output specific basis functions, such as polynomial trends and periodic seasonalities. In this paper, we repurpose the N-BEATS architecture, shifting its application from forecasting to unsupervised signal decomposition, exploiting its structural priors to separate continuous base loads from discrete lighting pulses.
